\chapter{Introduzione e pattern matching esatto}
\label{cap:introduzione}

\section{Presentazione del corso}
\label{sec:presentazione}

\scan{1}
Il corso di Algoritmi Avanzati è tenuto dal prof.\ Alberto Policriti.

\subsection{Le tre parti del corso}

Il corso si articola in tre parti, a ciascuna delle quali corrisponde un libro di
riferimento: \emph{string matching}, \emph{algoritmi randomizzati} e
\emph{strutture dati compresse}.

\begin{figure}[H]
  \centering
  % corso-tre-parti -- mappa delle tre parti del corso (quaderno, pag. 1).
\begin{tikzpicture}[
    >={Stealth[length=4.5pt,width=3.5pt]},
    font=\small,
    ramo/.style={->,semithick,shorten >=1.5pt},
    voce/.style={anchor=west,inner sep=1pt},
    glossa/.style={anchor=west,inner sep=1pt,font=\footnotesize},
  ]

  % --- livello 0: la radice ------------------------------------------------
  \node[anchor=west,inner sep=1pt] (radice) at (0,0) {3 parti};
  \coordinate (fork) at ($(radice.east)+(0.20,0)$);

  % --- livello 1: le tre parti del corso -----------------------------------
  \node[voce] (r1) at (2.00, 1.10) {\underline{string matching}};
  \node[voce] (r2) at (2.00, 0.00) {\underline{algoritmi randomizzati}};
  \node[voce] (r3) at (2.00,-1.10) {\underline{strutture dati compresse}};

  \draw[ramo] (fork) -- (r1.west);
  \draw[ramo] (fork) -- (r2.west);
  \draw[ramo] (fork) -- (r3.west);

  % --- diramazione: le due classi di algoritmi randomizzati -----------------
  \coordinate (v) at ($(r2.east)+(0.22,0)$);
  \node[glossa] (s1) at ($(v)+(0.40, 0.50)$) {Monte Carlo --- non sempre corretti};
  \node[glossa] (s2) at ($(v)+(0.40,-0.50)$) {Las Vegas --- sempre corretti};
  \draw[semithick] (v) -- ($(s1.west)+(-0.08,0)$);
  \draw[semithick] (v) -- ($(s2.west)+(-0.08,0)$);

  % --- glossa sull'obiettivo delle strutture dati compresse -----------------
  \node[glossa] (tilde) at ($(r3.east)+(0.14,0)$) {$\sim$};
  \node[glossa] (g3) at ($(tilde.east)+(0.08,0)$)
        {comprimere dati anche molto complessi};

\end{tikzpicture}

  \caption{Le tre parti in cui è suddiviso il corso, con le annotazioni a margine
  sulla dicotomia Monte Carlo / Las Vegas e sull'obiettivo delle strutture dati compresse.}
  \label{fig:corso-tre-parti}
\end{figure}

\begin{osservazione}
La distinzione fra i due tipi di algoritmi randomizzati riguarda dove viene «scaricata»
l'aleatorietà: un algoritmo \MonteCarlo ha tempo di esecuzione deterministico ma può
restituire una risposta errata (con probabilità piccola e controllabile), mentre un
algoritmo \LasVegas restituisce sempre la risposta corretta, e ad essere aleatorio è
invece il suo tempo di esecuzione.
\end{osservazione}

L'obiettivo della terza parte, infine, è comprimere dati anche molto complessi
mantenendoli interrogabili.

\subsection{Testi di riferimento e materiale}

\begin{enumerate}
  \item D.~Gusfield, \emph{Algorithms on Strings, Trees, and Sequences};
  \item R.~Motwani, P.~Raghavan, \emph{Randomized Algorithms};
  \item G.~Navarro, \emph{Compact Data Structures}.
\end{enumerate}

La numerazione dei testi ricalca quella delle tre parti del corso: a ogni modulo
corrisponde un testo canonico distinto. A questi si aggiungono gli appunti
(non ufficiali) del docente e i nostri appunti presi a lezione. Tutto il materiale
del corso è reso disponibile su Teams.

\subsection{Modalità d'esame}

L'esame è orale. Durante il corso vengono proposti spunti per approfondimenti;
l'esame consiste quindi in un (eventuale) approfondimento seguito dal colloquio orale.

\section{Il problema del pattern matching esatto}
\label{sec:pattern-matching-esatto}

\scan{2}
La prima delle tre parti del corso, quella dedicata al \emph{string matching}, comincia
dal problema che ne costituisce il nucleo.

\begin{definizione}[Pattern matching esatto]\label{def:pattern-matching-esatto}
\emph{Input.} Un alfabeto $\S$ di cardinalità finita e due stringhe $P,T\in\S^{*}$, con
$\abs{P}=n$ e $\abs{T}=m$ (di norma $n\le m$).

\emph{Output.} L'insieme delle posizioni
\[
  \set{i \mid \sub{T}{i}{i+n-1}=P}.
\]
\end{definizione}

Si vogliono cioè \emph{tutte} le occorrenze di $P$ (il \emph{pattern}) all'interno di $T$
(il \emph{testo}), non soltanto la prima.

\begin{osservazione}
Le stringhe sono indicizzate a partire da $1$ e $\sub{S}{i}{j}$ denota il fattore
$S[i]S[i+1]\cdots S[j]$. Perché il confronto abbia senso deve essere $1\le i\le m-n+1$:
le posizioni di allineamento del pattern sul testo sono dunque $m-n+1$.
\end{osservazione}

\begin{figure}[H]
\centering
% pm-allineamento -- il pattern P appoggiato sul testo T (quaderno, pag. 2).
\begin{tikzpicture}[
    >={Stealth[length=4pt,width=3pt]},
    stringa/.style={line width=1pt},
    car/.style={anchor=base,inner sep=0pt,font=\small},
  ]

  % --- il testo T ----------------------------------------------------------
  \draw[stringa] (0,0) -- (8.4,0);
  \node[car] at (0.30, 0.16) {$a$};
  \node[car] at (0.90, 0.16) {$b$};
  \node[car] at (1.50, 0.16) {$c$};
  \node[car] at (2.15, 0.16) {$\cdots$};
  \node[anchor=west,font=\small] at (8.6,0) {$T$};

  % --- il pattern P --------------------------------------------------------
  \draw[stringa] (0,-1.3) -- (3.4,-1.3);
  \node[car] at (0.30,-1.14) {$a$};
  \node[car] at (0.90,-1.14) {$b$};
  \node[car] at (1.50,-1.14) {$c$};
  \node[car] at (2.15,-1.14) {$\cdots$};
  \node[anchor=west,font=\small] (etP) at (3.6,-1.3) {$P$};

  % --- posizione corrente di confronto: stessa ascissa in T e in P ----------
  \draw[->] (0.30,-0.52) -- (0.30,-0.10);
  \draw[->] (0.30,-1.82) -- (0.30,-1.40);

\end{tikzpicture}

\caption{Il pattern $P$ allineato sul testo $T$: le due stringhe vengono confrontate
carattere per carattere a partire dalla posizione corrente (le frecce), e poi $P$ scorre
verso destra di una posizione alla volta.}
\label{fig:pm-allineamento}
\end{figure}

\section{L'algoritmo naive}
\label{sec:naive}

L'idea è provare tutti gli allineamenti, uno per uno, con due loop innestati: quello
esterno sceglie la posizione $i$ in cui il pattern viene appoggiato sul testo, quello
interno confronta $P$ con il fattore $\sub{T}{i}{i+n-1}$ carattere per carattere.

\begin{algorithm}[H]
\caption{Algoritmo naive per il pattern matching esatto}
\label{alg:naive}
\begin{algorithmic}[1]
\Require $T$ con $\abs{T}=m$, $P$ con $\abs{P}=n$
\Ensure tutte le posizioni $i$ tali che $\sub{T}{i}{i+n-1}=P$
\For{$i \gets 1$ \textbf{to} $m-n+1$}
    \Comment{allineamenti possibili}
    \State $j \gets 1$
    \While{$j \le n$ \textbf{and} $T[i+j-1] = P[j]$}
        \Comment{\textsc{compare}$(P,\sub{T}{i}{i+n-1})$}
        \State $j \gets j+1$
    \EndWhile
    \If{$j = n+1$}
        \State \textbf{output} $i$
        \Comment{$n$ confronti positivi: occorrenza in $i$}
    \EndIf
\EndFor
\end{algorithmic}
\end{algorithm}

Il ciclo interno --- la procedura \textsc{compare} --- scorre da sinistra a destra: o si
arresta al primo mismatch, e allora in posizione $i$ non c'è un'occorrenza, oppure arriva
in fondo dopo $\abs{P}=n$ confronti tutti positivi, e in tal caso $i$ viene emesso in
output. In entrambi i casi il pattern viene poi fatto scorrere di una sola posizione verso
destra: nessuna informazione raccolta durante i confronti viene riutilizzata.

\subsection{Complessità}
\label{sec:naive-complessita}

Il ciclo esterno esegue $m-n+1$ iterazioni e ogni chiamata a \textsc{compare} costa al più
$n$ confronti, da cui
\[
  \Oh(m\cdot n).
\]

Il limite è stretto (\emph{tight}): basta considerare il caso
\[
  P=a^{\,n-1}b, \qquad T=a^{\,m},
\]
cioè
\[
  T=\overbrace{a\,a\,a\,\cdots\,a\,a}^{m\ \text{volte}},
  \qquad
  P=\underbrace{a\,a\,\cdots\,a}_{n-1}\,b .
\]
Per ogni allineamento i primi $n-1$ confronti hanno successo e l'$n$-esimo fallisce: si
spendono esattamente $n(m-n+1)$ confronti --- cioè $\Th(m^{2})$ quando $n$ e $m-n$ sono
entrambi $\Th(m)$, con il massimo per $n\approx m/2$ --- e ciò nonostante l'output è vuoto.

D'altra parte non si può scendere sotto il costo della lettura dell'input. Per esempio, con
$n=1$ il naive costa $\Oh(m\cdot 1)=\Oh(m)$ ed è già ottimo: qualunque algoritmo deve almeno
leggere il tutto, e dunque
\[
  \Om(n+m).
\]

Per questo motivo la complessità $\Oh(m\cdot n)$ dell'algoritmo naive è detta
\emph{quadratica}. Fra questo upper bound e il lower bound banale $\Om(n+m)$ resta un
divario: la domanda naturale --- esiste un algoritmo sub-quadratico? --- è il punto di
partenza del capitolo successivo.
